\cotefile{16}
\scanpage{31}
\begin{flushright}
\unsure{Algèbres|Algèbre} de Hodge
\end{flushright}
% archive sheet number: 30

\scanpage{32}
\begin{flushright}
Formulaire des $L,\Lambda$, etc
\end{flushright}
% archive sheet number: 31

\scanpage{33}
\uline{Formulaire de} $L,\ \handwritten{\Lambda}$.

$X$ variété projective lisse connexe polarisée par $\handwritten{\unsure{\xi|\zeta}}_X\handwritten{\in} H^{\handwritten{2}}(X)$, $L_X$ opération de cup-produit associé\handwritten{e,} $\handwritten{\Lambda}_X$ opération de degré $-2$ correspondante, $p_i^X$ les projections de $H^{\handwritten{*}}(X)$ sur les morceaux primitifs $P^i(X)$ pour $0\handwritten{\leq}i\handwritten{\leq}2n$, $Y$ une section hyperplane \deleted{\unsure{d}}lisse de $X$, $\handwritten{\xi}_Y\handwritten{\in}H^{\handwritten{2}}(Y)$ la self-intersection (induite par $\handwritten{\xi}_X$), d'où $L_Y,\handwritten{\Lambda}_Y,p_Y^i$ ; $\phi:Y\handwritten{\to}X$ l'immersion, d'où $\phi^{\handwritten{*}}$ et $\phi_{\handwritten{*}}$.

A) \qquad On a
\[
\begin{aligned}
1)\quad &\operatorname{Im}L_X=\mathop{\handwritten{\bigcap}}_{0\handwritten{\leq}i\handwritten{\leq}n}\operatorname{Ker}p_{\deleted{\unsure{X}}i}^{\deleted{\unsure{i}}X},&
 \operatorname{Ker}L_X&=\mathop{\handwritten{\sum}}_{n\handwritten{\leq}i\handwritten{\leq}2n}\operatorname{Im}p_i^X,\\
2)\quad &\operatorname{Im}\handwritten{\Lambda}_X=\mathop{\handwritten{\bigcap}}_{n\handwritten{\leq}i\handwritten{\leq}2n}\operatorname{Ker}p_i^X&
 \operatorname{Ker}\handwritten{\Lambda}_X&=\mathop{\handwritten{\sum}}_{0\handwritten{\leq}i\handwritten{\leq}n}\operatorname{Im}p_i^X,
\end{aligned}
\]
d'où en particulier les formules
\[
\begin{aligned}
1')\quad &p_i^X L_X=0\quad\text{si}\quad 0\handwritten{\leq}i\handwritten{\leq}n,\quad&
 \deleted{\unsure{p}}L_Xp_i^X&=0\quad\text{si}\quad n\handwritten{\leq}i\handwritten{\leq}2n,\\
2')\quad &\deleted{\unsure{\phi}}p_{\deleted{\unsure{X}}i}^X\handwritten{\Lambda}_X=0\quad\text{si}\quad n\handwritten{\leq}i\handwritten{\leq}2n,\quad&
 \handwritten{\Lambda_X}p_i^X&=0\quad\text{si}\quad 0\handwritten{\leq}i\handwritten{\leq}n,
\end{aligned}
\]
et en particulier
\[
p_n^X L_X=L_Xp_n^X=p_n^X\handwritten{\Lambda}_X=\handwritten{\Lambda}_X p_n^X=0\quad. \tag*{3)}
\]

B) \qquad On a
\[
\handwritten{\Lambda}_X L_X=\mathrm{id}_X-\mathop{\handwritten{\sum}}_{\deleted{0}\handwritten{n}\handwritten{\leq}i\handwritten{\leq}\deleted{n}\handwritten{2n}}p_i^X,\qquad
 L_X\handwritten{\Lambda}_X=\mathrm{id}_X-\mathop{\handwritten{\sum}}_{\deleted{\unsure{n}}\handwritten{0}\handwritten{\leq}i\handwritten{\leq}\deleted{\unsure{2n}}\handwritten{n}}p_i^X, \tag*{4)}
\]
\deleted{d'où compte tenu de $1')$ et $2')$ les formules
\[
\begin{aligned}
4')\quad &\handwritten{\Lambda}_X L_X^2\handwritten{=}
 \deleted{\unsure{L_X\Lambda_X L_X}}\handwritten{\inserted{\left(\mathrm{id}_X-\sum_{0\leq i\leq 2n}p_i^X\right)L_X}},\quad
 \deleted{\text{et }\unsure{L_X\Lambda_X^2=\Lambda_XL_X\Lambda_X}}\handwritten{\inserted{L_X^2\Lambda_X=L_X\left(\mathrm{id}_X-\sum_{0\leq i\leq 2n}p_i^X\right)}},\\
4'')\quad & L_X\handwritten{\Lambda}_X L_X=L_X\handwritten{\Lambda}_X L_X\deleted{\unsure{\sum_{0\leq i\leq n}p_i^X=L_X\sum_{n\leq i\leq 2n}p_i^X}}\quad\handwritten{\text{et formules ``transposées''.}}
\end{aligned}
\]}

C) \qquad On a
\[
\phi_{\handwritten{*}}\handwritten{\Lambda}_Y\phi^{\handwritten{*}}=\mathrm{id}_X-\mathop{\handwritten{\sum}}_{0\handwritten{\leq}i\handwritten{\leq}2n}p_i^X
 \handwritten{\inserted{\overset{\text{par }4)}{=}(\Lambda_X L_X+L_X\Lambda_X)-\mathrm{id}_X+p_n^X}} \tag*{5)}
\]
d'où\handwritten{,} \deleted{compte tenu de $4'')$, $1')$ et $2')$}
\handwritten{\inserted{posant $(6)\ T_X=\mathrm{id}_X+\varphi_*\Lambda_Y\varphi^*$,}}
\deleted{\[
 5')\quad \handwritten{\Lambda}_X L_X^2=\deleted{\unsure{L_X\phi_*\Lambda_Y\phi^*}}\handwritten{\inserted{\varphi_*\Lambda_Y\varphi^* L_X}},\quad
 \handwritten{L_X^2\Lambda_X=L_X\varphi_*\Lambda_Y\varphi^*,}
\]}
\handwritten{\inserted{(donc $\Lambda_X L_X+L_X\Lambda_X=T_X-p_n^X$),}}

d'où compte tenu de $\deleted{\unsure{5'|3}}\handwritten{3\unsure{*|'}})$ la formule
\[
\begin{aligned}
\deleted{\unsure{6}}\handwritten{7})\quad
 L_X^2\handwritten{\Lambda}_X+2L_X\handwritten{\Lambda}_{\deleted{\unsure{Y|X}}X}L_X+
 \handwritten{\Lambda}_{\deleted{\unsure{Y|X}}X}L_X^2
 &=\deleted{L_X(\mathrm{id}_X\unsure{+|-}\phi_*\handwritten{\Lambda}_Y\phi^*)\unsure{+2|+}( \mathrm{id}_X\unsure{+|-}\phi_*\handwritten{\Lambda}_Y\phi^*)L_X}\\
 &\handwritten{\deleted{\unsure{2L_X+L_X}}\,\inserted{L_XT_X+T_XL_X,\quad\text{avec } T_X}}\\
 &\handwritten{\deleted{T_X=\mathrm{id}_X+\varphi_*\Lambda_Y\varphi^*}\quad\inserted{\text{donné par }(6)}}
\end{aligned}
\]

\uline{Corollaire} \quad $(A(X\mathbin{\mathrm{x}}X)$ et $C(Y))$ implique $C(X)$.

\handwritten{On y peut remplacer $A(X\times X)$ par $A^{\unsure{\circ|0}}(X\times X)$ (condition ``des \unsure{dim|dims|dim.} critiques'' $H^{\unsure{(2n)|2n|(2n+1)}-2}\to H^{(2n)+2}$}
% archive sheet number: 32

\scanpage{34}
\uline{Prop.}
\[
\left\{\begin{aligned}
 P^i(X)&\xrightarrow{\sim}P^i(Y)&&\text{si } i\leq n-2\\
 P^{n-1}(X)&\hookrightarrow P^{n-1}(Y)&&\deleted{\unsure{\text{si }i|\text{si }i=n-1}}\\
 \unsure{P^i|\mathrm{prim}}(X)&=\operatorname{Ker}\varphi^i&&\text{si }i\geq n
\end{aligned}\right.
\]
\diagram{\begin{tikzcd}[ampersand replacement=\&,column sep=large,row sep=large]
\deleted{\unsure{P|H}}H^i(X) \arrow[r,"{L_X^{n-i+1}}"] \arrow[d,"{\varphi^i}"',"{\unsure{\sim|2|\simeq}}"] \& H^{2n-i+2}(X) \& {}\\
H^i(Y) \arrow[r,"{L_Y^{n-i\deleted{\unsure{-1}}}}"] \& H^{2(n-1)\deleted{\unsure{+}}-i}(Y) \arrow[u,"{\varphi_{\unsure{2(n-1)+i|2(n-1)-i}}}"',"{\sim}"] \& \smash{\raisebox{2.2ex}{$\left.\rule{0pt}{2.8em}\right]\ \deleted{\unsure{\text{si}}}\ i\leq n-2$}}
\end{tikzcd}}
\[
i\leq n-2\Rightarrow 2(n-1)-i\geq n
\]
\diagram{\begin{tikzcd}[ampersand replacement=\&]
H^j(X)\arrow[d,"{\varphi^*}"]\arrow[dr,"{L_X}"]\& \&\\
H^j(Y)\arrow[r,"\sim"]\& H^{j+2}(X)\& j\geq n
\end{tikzcd}}

\uline{\deleted{\unsure{Cor}}}

\unsure{Notn|Notons|Non} $\longrightarrow$
\[
\left\{\begin{aligned}
H^*(X)&=\deleted{\unsure{\sum_i}}\coprod_{\substack{i+j\leq n\\i,j\geq0}}L_X^jP^i(X)
=\coprod_{\substack{i+j\leq n-1\\i,j\geq0}}L_X^jP^i(X)\ \oplus\coprod_{j\geq n}P^j(X)\\
H^*(Y)&=\coprod_{\substack{i+j\leq n-1\\i,j\geq0}}L_Y^jP^i(Y)
\end{aligned}\right.
\]

\uline{Cor} \quad $\varphi^*$ induit un mor d'objets \unsure{bigradués}
\[
\deleted{\unsure{-}}\coprod_{\left\{\substack{i+j\leq n-1\\i,j\geq0}\right.} L_X^j P^i(X)
 \longrightarrow H^*(Y)=\deleted{\unsure{+}}\coprod_{\substack{i+j\leq n-1\\i,j\geq0}} L_Y^jP^i(Y)\ ; 
\]
\deleted{\unsure{qui est un}} \inserted{\unsure{C'est|Ceci}} induit un \unsure{isom} sur les morceaux bihomogènes, sauf $L_X^0P^{n-1}(X)=P^{n-1}(X)\longrightarrow\deleted{\unsure{P^{n-2}}}\,L_Y^0P^{n-1}(Y)=P^{n-1}(Y)$.

\uline{Cor} \quad Le noyau de $\varphi^*$ est \inserted{égal à celui de $L_X$, i.e.} $\deleted{\unsure{\sum}}\coprod_{j\geq n}P^i(X)$ \deleted{(qui est \unsure{un|aussi un} idéal)}. Le conoyau est
\[
\begin{aligned}
E^{n-1}&=H^{n-1}(Y)/\operatorname{Im}H^{n-1}(X)=P^{n-1}(Y)/\operatorname{Im}P^{n-1}(X)\\
 &\simeq\operatorname{Ker}(\varphi_{\unsure{n-1|*}}\colon H_{\unsure{n-1}}^{n-1}(Y)\unsure{-|\to}H^{n-1}(X))
\end{aligned}
\]
% archive sheet number: 33

\scanpage{35}
\mspage{\handwritten{I}}
\margin{\handwritten{$\uparrow$}}
\begin{center}
Cohomologie des catégories de Banach

\rule{8em}{0.4pt}

par Max Karoubi
\end{center}

Soit $\handwritten{\mathscr{C}}$ une catégorie additive \handwritten{\textsuperscript{$(*)$}}. Le \uline{groupe de Grothendieck} $K(\handwritten{\mathscr{C}})$ de la catégorie $\handwritten{\mathscr{C}}$ est le quotient du groupe libre engendré par les classes d'isomorphie $\handwritten{\{}G\handwritten{\}}$ d'objets $G$ de $\handwritten{\mathscr{C}}$ par le sous-groupe engendré par les relations $\handwritten{\{}E+F\handwritten{\}}-\handwritten{\{}E\handwritten{\}}-\handwritten{\{}F\handwritten{\}}=0$. Pour des applications topologiques du groupe de Grothendieck ,voir par exemple \textsuperscript{(1)}. Dans cette note on se propose de définir de même ,pour des catégories additives munies d'une certaine structure ,des groupes $K^n(\handwritten{\mathscr{C}})$ , $n\handwritten{\in}\unsure{Z|\mathbb{Z}}$ ,avec $K^0(\handwritten{\mathscr{C}})\handwritten{\approx}K(\handwritten{\mathscr{C}})$ ,jouissant de certaines propriétés ``cohomologiques'' en relation avec les théorèmes de périodicité de Bott .

I \qquad \uline{Définitions}

1.1 \quad Une \uline{catégorie de Banach} (sous-entendu sur $k$ avec $k=\mathbb{R}$ ou $\mathbb{C}$ ) est la donnée d'une catégorie avec objet nul $\handwritten{\mathscr{C}}$ et , quels que soient les objets $M$ et $N$ de $\handwritten{\mathscr{C}}$ ,d'une structure de $k$-espace de Banach sur l'ensemble $\operatorname{Hom}_{\handwritten{\mathscr{C}}}(M,N)$. Cette donnée est astreinte aux \deleted{\unsure{c|C}}conditions suivantes :

a) Quels que soient les objets $M$ ,$N$ et $P$ de $\handwritten{\mathscr{C}}$ ,l'application
\[
\operatorname{Hom}_{\handwritten{\mathscr{C}}}(M,N)\mathbin{\mathrm{x}}\operatorname{Hom}_{\handwritten{\mathscr{C}}}(N,P)\longrightarrow\operatorname{Hom}_{\handwritten{\mathscr{C}}}(M,P)
\]
définie par la composition des morphismes ,est bilinéaire

\noindent\rule{\linewidth}{0.4pt}

\handwritten{\textsuperscript{$(*)$}} Toutes les catégories considérées ici sont des $\handwritten{\mathscr{U}}$-catégories au sens de Grothendi\handwritten{e}ck \textsuperscript{(3)}. Les objets d'une catégorie forment donc en particulier un ensemble .
% archive sheet number: 34

\scanpage{36}
\unsure{On aura donc additivement}
\[
 H^*(Y)=\coprod_{\substack{i+j\leq n-1\\i,j\geq0}}L_X^jP^i(X)\quad\oplus\quad E^{n-1}\tag*{\unsure{$(\alpha)$|$(b)$}}
\]

L'opération
\[
\varphi_*:H^*(Y)\longrightarrow H^*(X)
\]
s'identifie \unsure{alors|donc} à $L_X$ sur le morceau $\coprod_{\substack{i+j\leq n-1\\i,j\geq0}}L_X^jP^{\unsure{j|i}}(X)$ \unsure{du premier morceau de|des premiers morceaux de} $(\alpha)$, (et est \deleted{\unsure{nul sur}} donc injectif sur ce morceau) et il est \uline{nul} sur le morceau $\deleted{\unsure{\coprod_{0\leq i\leq n-1} L_X^{n-1-i}P^i(X)\oplus|\coprod_{0\leq i\leq n-1}L_X^jP^i(X)\oplus}}\ E^{n-1}$, \deleted{qui correspond \unsure{par l'isom} \illegible{1} à $\sum_{n-1\leq i\leq2(n-1)}P^i(Y)$ \unsure{$*$}}. Ainsi on trouve le

\uline{Corollaire}
\[
\begin{aligned}
\operatorname{Ker}\varphi_* &\deleted{\unsure{=\sum_{n-1\leq j\leq2n-1}P^j(Y)=\operatorname{Ker}}}\ E^{n-1}(Y)\\
\operatorname{Im}\varphi_* &=\operatorname{Im} L_X=\bigcap_{0\leq i\leq n}\operatorname{Ker}p_i^{\unsure{X|\alpha}}
\end{aligned}
\]
(la dernière formule, car $H^*(Y)=\operatorname{Im}\varphi^*+E^{n-1}$ et $\varphi_*$ s'annule sur $E^{n-1}$, \unsure{donc} $\operatorname{Im}\varphi_*=\operatorname{Im}\varphi_*\varphi^*\deleted{\unsure{=I}}$, or $\varphi_*\varphi^*=L_X$).
\[
\left\{\begin{aligned}
p_i^X\varphi_*&=0&&\text{si }0\leq i\leq n\\
p_j^X\varphi_*&=\varphi_*p_{j-2}^Y&&\text{si }\unsure{n+1|\deleted{n+1}}\leq i\leq 2n.
\end{aligned}\right.
\]

\margin{\deleted{\unsure{fonctorialité|variation}}\\\unsure{variantes}}
On a
\[
\begin{aligned}
\left\{\begin{aligned}
\varphi^*L_X&=L_Y\varphi^*\\
L_X\varphi_*&=\varphi_*L_Y
\end{aligned}\right. \qquad&
\begin{aligned}
\varphi_*\varphi^*&=L_X,\quad\varphi^*\varphi_*=L_Y\\
&\text{\unsure{formule des projections|formule de projection|formules de projection}}
\end{aligned}
\end{aligned}
\]
\deleted{Pour la \unsure{deuxième|dernière} formule, on \unsure{note} que les deux \unsure{mbs} sont nuls sur $E^{n-1}$, donc il suffit de vérifier,
\[
 (L_X\varphi_*)\varphi^*=(\varphi_*L_Y)\varphi^*\quad\text{i.e.}\quad
 L_X(\varphi_*\varphi^*)=\varphi_*\underbrace{(L_Y\varphi^*)}_{\inserted{\varphi_*\varphi^*L_X}}=
\]
or les deux mbs sont $L_X^2$.}
% archive sheet number: 35

\scanpage{37}
\mspage{-8-}
Now we apply these results to locally noetherian preschemes, giving first a special case.

\uline{Lemma} \uline{3.3}: Let $A$ be a noetherian ring, let $\unsure{\mathscr{P}|\mathfrak{p}}$ be a prime ideal of $A$, let $I$ be an injective hull of $k(\mathscr{P})$, and let $X=\operatorname{Spec} A$. Then the sheaf $\widetilde{I}$ is injective in the category $C_m$ of all $\mathscr{O}_X$-modules.

\uline{Proof}: Let $x\in X$ be the point corresponding to the prime ideal $\mathscr{P}$. Then $\widetilde{I}$ is the direct limit of the sheaves $\widetilde{E}_i$, where the $E_i$ are as in Proposition 3.2. Since $E_i$ is an Artin module over the local ring $A_{\mathscr{P}}$, $\widetilde{E}_i$ is the constant sheaf $E_i$ on $\overline{\{X\}}$, and $0$ elsewhere. Thus $\widetilde{I}$ is the constant sheaf $I$ on $\overline{\{X\}}$, and $0$ elsewhere.

We now show directly that $\widetilde{I}$ is an injective $\mathscr{O}_X$-module. Indeed, if $F$ is any $\mathscr{O}_X$-module, to give a sheaf homomorphism from $F$ into $\widetilde{I}$ is the same thing as to give a homomorphism from the stalk $F_x$ into $I$ ! Thus if $F\subseteq G$, and we are given a homomorphism $F\longrightarrow\widetilde{I}$, it extends to $G$, since we need only extend the map to the stalk $G_x\supseteq F_x$ , and $I$ is an injective $\mathscr{O}_{x,X}$-module.

\uline{Proposition} \uline{3.4}: Let $X$ be a locally noetherian prescheme, and let $F$ be a quasi-coherent $\mathscr{O}_X$-module. Then the following conditions are equivalent:

(i) $F$ is injective in the category $C_m$ of all $\mathscr{O}_X$-modules,

(ii) $F$ can be written as a direct sum $F=\oplus F_\alpha$ , where each $F_\alpha$ is the constant sheaf $I_x$ on $\overline{\{x\}}$ and $0$ outside, for some $x\in X$, and where $I_x$ is an injective hull of the residue field
% archive sheet number: 36

\scanpage{38}
\deleted{\unsure{A|A)}} Moyennant l'isom $(\alpha)$, l'opération $L_Y$ \deleted{\unsure{s'id}}s'identifie à $L_X$ sur le morceau $\coprod_{\substack{i+j\leq n-2\\i,j\geq0}}L_X^jP^i(X)$, et son noyau est le morceau $\coprod_{\substack{\deleted{\unsure{i+j}}\\i+j=n-1\\i,j\geq0}}\underline{L_X^jP^{\unsure{j|i}}(X)}\oplus E^{(n-1)}=\coprod_{i\geq n+1}P^j(X)\oplus E^{(n-1)}$ (qui n'est d'ailleurs autre que $\coprod_{\substack{\deleted{\unsure{i\geq}}\\n\leq i\leq2(n-1)}}P^i(Y)$, comme il se doit). L'opération $\Lambda_Y$ est nulle sur $E^{n-1}$, et coïncide sur le morceau $\coprod_{\substack{i+j\unsure{=|\leq} n-1\\i,j\geq0}}L_X^jP^i(X)$ avec $\Lambda_X$. On retrouve \deleted{\unsure{l'égalité|l'injectivité}} \inserted{que le noyau de} $\Lambda_Y$ \deleted{\unsure{sur}} est $\coprod_{0\leq i\deleted{\unsure{+j}}\leq n-1}P^i(X)+E^{n-1}=\coprod_{0\leq i\leq \unsure{n-1|n+1}}P^i(Y)$, \deleted{\unsure{comme}} et que $\Lambda_Y$ est injectif sur le reste $\deleted{\unsure{=|\longrightarrow}}\coprod_{\substack{i+j\leq n-2\\\deleted{\unsure{i\geq0,\ j\geq0}}\\i,j\geq0}}L_X^{1+j}P^i(X)$.

On en déduit
\deleted{\[
\left\{\begin{aligned}
 \deleted{\unsure{\Lambda_Y\varphi^*=\Lambda_X}}\,\varphi^*\Lambda_X&=\Lambda_Y\varphi^*+\\
 \varphi_*\Lambda_Y&=\Lambda_X\varphi_*
\end{aligned}\right.
\]
Pour la \deleted{deuxième}\inserted{\unsure{pre}} relation, on la vérifie \unsure{d'abord} sur les $P^j(X)\unsure{=}$ pour $j\geq n$, \inserted{i.e. les $L_X^{n-i}P^i(X)$ pour $i\leq n-1$ ;} \deleted{le premier \unsure{mb}} y est nul car $\varphi^*$ \unsure{l'est}, le deuxième}

\deleted{$\unsure{\varphi\Lambda_X=}$}
\margin{$\Lambda_Y^m\varphi^*=\varphi^*\Lambda_X^m(\mathrm{id}-\sum)$\\pour \unsure{tt} $m$ $\Longleftarrow$}
\[
\left\{\begin{aligned}
 \Lambda_Y\varphi^*&=\varphi^*\Lambda_X\left(\mathrm{id}_X-\sum_{n+1\leq j\leq2n}p_X^j\right)
 =\left(\mathrm{id}_Y-\unsure{\bigl(}\sum_{\deleted{\unsure{+}}n-1\leq j\leq 2(n-1)}p_X^j\right)\varphi^*\Lambda_X\\
 \varphi^*\Lambda_Xp_X^j&=p_Y^{j-2}\varphi^*\Lambda_X\deleted{\unsure{=p_Y^{j-2}\Lambda_Y\varphi^*}}\qquad\text{si }j\geq n+1
\end{aligned}\right.
\]
\[
\left\{\begin{aligned}
 \deleted{\unsure{\Lambda}}\varphi_*\Lambda_Y&=\left(\mathrm{id}_X-\sum_{0\leq i\leq n-1}p_X^i\right)\Lambda_X\varphi_*=\Lambda_X\varphi_*\left(\mathrm{id}_Y-\sum_{0\leq i\leq n-1}p_Y^i\right)\\
 p_X^i\Lambda_X\varphi^*&=\Lambda_X\varphi^*p_Y^i=p_X^i\varphi_*\Lambda_Y\qquad\text{si }0\leq i\leq n-1
\end{aligned}\right.
\]
% archive sheet number: 37

\scanpage{39}
\mspage{-7-}
and such that if $N$ is any non-zero submodule of $I$, then $M\cap N\neq(0)$.

Eckmann and Schopf have shown that every module has an injective hull, unique up to non-unique isomorphism [``Uber injektive Moduln'', Archiv. der. Math. \uline{4} (1953)]. We recall the following theorem, due to Matlis [``Injective modules over noetherian rings'', Pac. J. Math. \uline{8} (1958) 511-528]:

\uline{Theorem} \uline{3.1}: Let $A$ be a noetherian ring, and let $I$ be an injective $A$-module. Then $I$ can be written uniquely as a direct sum $I=\oplus I_\alpha$, where each $I_\alpha$ is the injective hull of the residue field $k(\unsure{\mathscr{P}|\mathfrak{p}}_\alpha)$ of a suitable prime ideal $\mathscr{P}_\alpha$ of $A$.

We also have information about the structure of one of these injective hulls [Matlis, ibid ]:

\uline{Proposition} \uline{3.2}: Let $A$ be a noetherian ring, let $\mathscr{P}$ be a prime ideal of $A$, and let $I$ be an injective hull of the residue field $k(\mathscr{P})$ of $\mathscr{P}$. Then one can write $I$ as a direct limit of submodules
\[
0\subseteq E^1\subseteq E^2\subseteq\ldots
\]
where each $E^i$ is an Artin module over the local ring $A_{\mathscr{P}}$, and where, for each $i$,
\[
\dim_{k(\mathscr{P})}(E^{i+1}/E^i)=\dim_{k(\mathscr{P})}\left((\mathscr{P}^i/\mathscr{P}^{i+1})\otimes_A k(\mathscr{P})\right).
\]
% archive sheet number: 38

\scanpage{40}
On conclut du premier groupe
\[
\varphi_*\Lambda_Y\varphi^*
 =\underbrace{\varphi_*\varphi^*\Lambda_X}_{\substack{L_X\Lambda_X\\ \shortparallel\\ \mathrm{id}_X-\sum_{0\leq i\leq\unsure{n|2n}}p_i^X}}
 \left(\mathrm{id}_X-\sum_{n+1\leq j\leq2n}p_X^j\right)
 =\underbrace{\mathrm{id}_X-\sum_{0\leq j\leq2n}p_X^j}_{q}.
\]
$[$Plus gén.
\[
\varphi_*\Lambda_Y^m\varphi^*
=\deleted{\unsure{\sum|\Sigma}}\left(\mathrm{id}-\sum_{0\leq i\leq n}p_X^i\right)\Lambda_X^{m-1}
 \left(\mathrm{id}-\sum_{n+1\leq j\leq2n}p_X^j\right)=q\Lambda_X^{m-1}q
\]
où $q=\mathrm{id}-\sum_{0\leq i\leq2n}p_X^i$

Aussi :
\[
\varphi^*\Lambda_X\varphi_*=\mathrm{id}_Y-p_E
\]
\inserted{projection sur partie évanescente $E^{n-1}$}

d'où, en multipliant par $\Lambda_Y$ à gauche, comme $\Lambda_Yp_E=0$
\[
\Lambda_Y=\underline{\Lambda_Y\varphi^*}\Lambda_X\varphi_*=
 \underline{\varphi^*\Lambda_X\Bigl(\mathrm{id}_X-\sum_{n+1\leq j\leq2n}}
 \overbrace{\underline{p_X^j\Bigr)}\,\Lambda_X}^{\substack{0\\\shortparallel}}\varphi_*\qquad\text{donc}
\]
\[
\boxed{\Lambda_Y=\varphi^*\Lambda_X^2\varphi_*}
\qquad\text{Donc}\quad C(X)\quad\text{\deleted{\unsure{équivaut|implique}}\inserted{implique}}\quad C(Y).
\]

Soit $p_X^{i,j}$ \unsure{la|les} projection sur $L_X^{\unsure{j|t}}P^i(X)$ \quad $(i+j\leq n,\ i,j\geq0)$.

Alors
\margin{$\unsure{\lambda|L|\Lambda}$?}
\[
\varphi_*p_Y^{i,j}\inserted{\unsure{\lambda|\Lambda}}\varphi^*=p_X^{i,j+1}
 \qquad\text{pour}\qquad \unsure{0|i}+j\unsure{<|\leq}n-1,\quad i,j\geq0\quad\deleted{\unsure{\text{algébr}}}
\]
donc la condition $D(Y)$ implique que les $p_X^{ij}$ sont algébriques pour $\deleted{\unsure{j}}\,i+j<n$, $i\geq0$ et \uline{$j\geq1$}, i.e. les seuls \unsure{évt. non alg.|int. en susp.|rest. en susp.} sont les \uline{$p_X^i$ pour $0\leq i\leq\unsure{2n|n|\deleted{2n}}$}. Mais
% archive sheet number: 39

\scanpage{41}
\mspage{-5-}
is the identity functor. Again since the question is local on $X$, and $L$ is coherent, we write an exact sequence
\[
0\longrightarrow L'\longrightarrow\mathscr{O}_X^p\longrightarrow L\longrightarrow0,
\]
with $L'$ also coherent. This gives an exact sequence
\[
\underline{\mathrm{Hom}}(\mathscr{O}_X^p,M)\longrightarrow\underline{\mathrm{Hom}}(L',M)\longrightarrow\underline{\mathrm{Ext}}^1(L,M)\longrightarrow0,
\]
which shows that $\underline{\mathrm{Ext}}^1(L,M)$ is coherent (resp. quasi-coherent), using the case $i=0$ applied to $L'$. We also get isomorphisms for all $i>1$,
\[
\underline{\mathrm{Ext}}^i(L,M)\cong\underline{\mathrm{Ext}}^{i-1}(L',M),
\]
which, by induction on $i$, shows that $\underline{\mathrm{Ext}}^i(L,M)$ is coherent (resp. quasi-coherent) for all $i$.

\uline{Corollary} \uline{2.2}: Let $X$ be a locally noetherian prescheme. Let $L^\bullet\in D(X)$ have coherent cohomology, and let $M^\bullet\in D^+(X)$ have coherent (resp. quasi-coherent) cohomology. Assume either

a) $L^\bullet\in D^-(X)$ , \qquad or

b) $L^\bullet\in D^+(X)$, and there exists an $n_{\unsure{0|\mathrm{o}}}\in\mathbb{Z}$ such that for all $i>n_0$, and for all coherent sheaves $F$, $\underline{\mathrm{Ext}}^i(F,M^\bullet)=0$.

Then $\text{\uuline{R}}\ \underline{\mathrm{Hom}}^\bullet(L^\bullet,M^\bullet)$ has coherent (resp. quasi-coherent) cohomology.

\uline{Proof}: The proof is a good exercise in the use of way-out functors, and will be left to the reader. Note that in case b) the hypothesis on $M^\bullet$ is needed to ensure that the functor
% archive sheet number: 40

\scanpage{42}
Moyennant $C(Y)$, $C(X)$ équivaut à la condition :

\noindent\begin{tabular}{@{}|p{.92\linewidth}@{}}
$\forall i\leq n-1$, l'hom $\psi_\xi^i\colon H^i(Y)\longrightarrow H^i(X)$, inverse à gauche naturel de $\varphi^i\colon H^i(X)\unsure{\mathord{-}|\longrightarrow}H^i(Y)$, \deleted{\unsure{est}} égal à $\Lambda_X\varphi_{\unsure{i|*}}\mid H^i(Y)$, est induit par une classe algébrique.
\end{tabular}

ou encore, :

\noindent\begin{tabular}{@{}|p{.92\linewidth}@{}}
$\Lambda_X\varphi_*$ est algébrique
\end{tabular}

\deleted{ou encore} N.B. $\Lambda_X\varphi_*$ correspon\deleted{\unsure{d à|d}}d à une classe \inserted{dans $\sum H^{2(n-1)-i}(Y)\otimes H^i(X)=H^{2(n-1)}(X\times Y)=H^{m-1}(X\times Y)$, $m=\dim X\times Y$} de cohomologie sur $X\times Y$, que nous polarisons par $\xi_X\otimes1+1\otimes\xi_Y$. Or \deleted{\unsure{$\xi_X\otimes1+$|$\zeta_X\otimes1+$}} $L_{X\times Y}(\Lambda_X\varphi_*)$ correspond à
\[
\begin{aligned}
\underbrace{\Lambda_X\varphi_*L_Y}_{\substack{\deleted{\unsure{\Lambda_X\varphi_*\varphi^*}}\\ \Lambda_XL_X\varphi_*}}
+\underbrace{L_X\Lambda_X\varphi_*}_{\deleted{\unsure{\varphi_*\Lambda_Y\varphi^*\varphi_*}}}
&=\deleted{\left(2\,\mathrm{id}-\sum_{n\leq i\leq2n}p_i^X\right)\varphi_*}\\
&(\Lambda_X L_X+L_X\Lambda_X)\varphi_*\\
&=(\varphi_*\Lambda_Y\varphi^*+\mathrm{id}_X)\varphi_*
\end{aligned}
\]
(car $p_n^X\varphi_*=0$)

Donc c'est algébrique. Donc on trouve
\[
 C(X)\Longleftrightarrow\bigl[C(Y)\text{ et }A^\circ(X\times Y)\bigr]
\]
puis de proche en proche
\[
 C(X)\Longleftrightarrow\bigl(A^\circ(X\times Y)\text{ et }A^\circ(Y\times Z)\text{ et }\deleted{\unsure{A}}A^\circ(Z\times T)\;{-}{-}\bigr)
\]
\deleted{$X\times Y,\,Y\times\unsure{Z}$} $X\supset Y\supset Z\supset T\,{-}{-}$ est une suite décroissante de sections hyperplanes non singulières.

% stamp: repeated diagonal photographic-paper watermark, apparently Gevaert.
% archive sheet number: 41

\scanpage{43}
\mspage{- IV-\handwritten{977}- \handwritten{12 -}}

$u:B\to A$, où $B=k\handwritten{[[}T_1,..,T_r\handwritten{]]}$ ($r$ entier quelconque) faisant de $A$ une $B$-algèbre finie, $\underline{m}$ est le seul point de $\operatorname{Spec}(A)$ au-dessus'\deleted{\unsure{du|de}}de l'idéal maximal \deleted{\illegible{1}} $\underline{n}$ de $B$ ; \deleted{\illegible{1}} le morphisme \deleted{\unsure{$\operatorname{Spec}(A)\to\operatorname{Spec}(B)$}} \inserted{\handwritten{$\operatorname{Spec}(A)\to\operatorname{Spec}(B)$}} ne peut être non ramifié en $\underline{m}$ que si $A/\underline{n}A$ est un corps, extension séparable de $k$ (17.4.1), ce qui implique que le corps résiduel $K$ de $A$ est une extension séparable de $k$. Si cette condition n'est pas remplie, la conclusion de (18.11.5\deleted{\unsure{p|b}}) ne peut jamais être vérifiée pour $\underline{p}=\underline{m}$, quel que soit l'entier $m$.

(ii) Supposons toujours $\underline{p}\ne\underline{m}$, mais considérons le cas où la condition de (18.11.5) relative à \inserted{\handwritten{l'existence de}} l'idéal minimal $\underline{q}$ n'est pas remplie. Soit $\underline{J}$ le noyau de l'homomorphisme canonique $A\to A_{\underline{p}}$, autrement dit l'ensemble des annulateurs des éléments de $A-\underline{p}$, et posons $A'=A/\underline{J}$ ; on a nécessairement $\underline{J}\subset\underline{p}$, et si l'on pose $\underline{p}'=\underline{p}/\underline{J}$, l'homomorphisme canonique $A_{\underline{p}}\to A'_{\underline{p}'}$ est \uline{bijectif}. En vertu de (18.11.3.\unsure{3|2}) et de (\uline{0},20.4.12), l'homomorphisme canonique \deleted{\unsure{$\Omega^1_{A/k}\otimes_A A'$}} $\widehat{\Omega}^1_{A/k}\otimes_A A'\to\widehat{\Omega}^1_{A'/k}$ est \uline{surjectif}, donc il en est de même de $(\widehat{\Omega}^1_{A/k})_{\underline{p}}\to(\widehat{\Omega}^1_{A'/k})_{\underline{p}'}$. D'autre part, l'homomorphisme canonique $A'\to A'_{\underline{p}'}$ est par définition \uline{injectif}, et les idéaux \inserted{premiers} minimaux de $A'$ correspondent donc biunivoquement à ceux de $A'_{\underline{p}'}$ ;
% slanted
\margin{\handwritten{\uline{\unsure{Vider|Viduer}} (ii) sauf \unsure{raison} \unsure{sérieuse} !}}
\deleted{\unsure{donc}} \inserted{\handwritten{par suite}} il y a un idéal minimal $\underline{q}'$ de $A'$ \deleted{\illegible{1}} contenu dans \deleted{\illegible{1}} $\underline{p}'$ et tel que $\dim(A'/\underline{q}')=\dim(A')$. On peut ainsi appliquer (18.11.5) à $A'$ et $\underline{p}'$, et il y a donc un $k$-homomorphisme local $u':B\to A'$, avec $B=k\handwritten{[[}T_1,..,T_m\handwritten{]]}$ tel que le morphisme correspondant \deleted{\unsure{$\operatorname{Spec}$}} \inserted{$f':$} $\operatorname{Spec}(A')\to\operatorname{Spec}(B)$ soit fini et non ramifié au point $\underline{p}'$. Mais on peut factoriser $u'$en $B\xrightarrow{u}A\xrightarrow{p}A'$ où \deleted{\illegible{3}} $p$ est l'homomorphisme canonique ; d'autre part, l'injection canonique $\operatorname{Spec}(A')\to\operatorname{Spec}(A)$ est un\deleted{e} \deleted{\unsure{homomorphisme}} \uline{isomorphisme local} au point $\underline{p}'$ en vertu de l'isomorphisme $A_{\underline{p}}\xrightarrow{\sim}A'_{\underline{p}'}$ et de (\uline{I},6.5.4). On voit donc que l'hypothèse a) de (18.11.5) entraîne encore, dans ce cas, l'existence d'un $k$-homomorphisme local $u:B\to A$ tel que

% archive sheet number: 42

